%BAB V: HARDSKILL - PENGEMBANGAN BACKEND
\section{MBKM Hardskill: Pengembangan Backend}

\begin{table}[H]
    \centering
    \begin{tabular}{ll}
        \textbf{Nama} & : Muhammad Argya Vityasy \\
        \textbf{NIM} & : 23/522547/PA/22475 \\
        \textbf{Pembimbing Akademik} & : Dr.\ techn.\ Guntur Budi Herwanto, S.Kom., M.Cs. \\
        \textbf{Pembimbing Program MBKM} & : Dr.\ techn.\ Guntur Budi Herwanto, S.Kom., M.Cs. \\
        \textbf{Mata Kuliah} & : Internship: Pengembangan Backend \\
        \textbf{Kode MK} & : MII21-3014 \\
        \textbf{Total SKS} & : 4 SKS \\
        \textbf{Judul Magang} & : MBKM Mandiri GDP Labs: Digital Employee \\
        \textbf{Durasi} & : 9 Februari 2026 -- 22 Juni 2026 \\
    \end{tabular}
\end{table}

\subsection{Landasan Teori}

Backend Digital Employee menangani orkestrasi agen AI, autentikasi, manajemen data, dan komunikasi antar layanan. Identity and Access Management (IAM) menjadi komponen kritis karena setiap permintaan ke sistem harus tervalidasi dan sesuai peran (role) pengguna atau agen.

GDP Labs mengembangkan GL-IAM SDK, kit perangkat lunak yang menyediakan abstraksi SIMI (Single Interface Multiple Implementation) untuk manajemen identitas dan akses. SDK ini mendukung berbagai penyedia autentikasi (PostgreSQL, Keycloak, Stack Auth) melalui satu interface yang sama, sehingga aplikasi tidak perlu menulis ulang kode autentikasi ketika penyedia diganti.

\subsection{Kegiatan}

\subsubsection{Atomicitas \texttt{create\_user\_with\_password} di GL-IAM SDK}

Kontribusi utama saya pada backend adalah perbaikan atomicitas operasi \texttt{IAMGateway.create\_user\_with\_password()} di repositori \texttt{gl-sdk}. Implementasi awal memakai atomicitas semu: pembuatan pengguna dan pemasangan kata sandi dilakukan dalam transaksi terpisah yang dibungkus \texttt{try/except} dengan rollback manual. Jika \texttt{set\_user\_password()} gagal setelah baris pengguna terkomit, tersisa pengguna tanpa kredensial (orphan user) di basis data.

Perbaikan dilakukan dalam dua tahap:

\begin{enumerate}
    \item \textbf{Rollback pada gateway}: Jika \texttt{set\_user\_password()} gagal, pengguna yang baru dibuat langsung dihapus. Perilaku ini diuji melalui \texttt{test\_create\_user\_with\_password\_rollback\_on\_failure} dan \texttt{test\_create\_user\_with\_password\_rollback\_continues\_on\_error}.
    \item \textbf{Atomicitas di level penyedia}: Tanggung jawab atomicitas dipindahkan dari IAMGateway ke \texttt{PostgreSQLUserStoreMixin}. Metode \texttt{create\_user\_with\_password} ditambahkan ke protokol \texttt{UserStoreProvider} dengan fallback dua langkah secara default, sehingga penyedia tanpa kata sandi lokal tidak terdampak. \texttt{PostgreSQLUserStoreMixin} meng-override-nya dengan implementasi atomik satu sesi dan satu komit: baris pengguna, kredensial, external identity, dan peran dikomit dalam satu transaksi. IAMGateway cukup mendelegasikan ke penyedia, tanpa mekanisme rollback yang rumit.
\end{enumerate}

\begin{lstlisting}[caption={Implementasi atomic \texttt{create\_user\_with\_password} pada \texttt{PostgreSQLUserStoreMixin}}]
# PostgreSQLUserStoreMixin: single-session atomic
async def create_user_with_password(self, input: UserCreateInput):
    async with self.session() as session:
        user = User(**input.to_dict(exclude={"password"}))
        session.add(user)
        credential = Credential(user_id=user.id, ...)
        session.add(credential)
        if input.external_identity:
            ext = ExternalIdentity(user_id=user.id, ...)
            session.add(ext)
        await session.commit()
        return user
\end{lstlisting}

Perubahan pendukung lainnya: penambahan field \texttt{external\_identity} pada \texttt{UserCreateInput} (diakses via \texttt{to\_dict()} untuk dukungan JIT provisioning), penanganan \texttt{InvalidCredentialsError} di gateway (dipetakan ke \texttt{VALIDATION\_ERROR}), pembaruan 8 pengujian gateway, serta penambahan \texttt{pyright} untuk pemeriksaan tipe LSP. Perubahan ini dirangkum dalam PR \#3645.

\subsubsection{GL-IAM Cookbook: API Key Hierarchy dan Tutorial DPoP}

Pada repositori \texttt{gl-iam-cookbook}, saya berkontribusi pada beberapa contoh integrasi:

\begin{itemize}
    \item \textbf{API Key Hierarchy}: Menyempurnakan contoh hierarki API key tiga tingkat (PLATFORM, ORGANIZATION, PERSONAL) di mana key tingkat atas dapat menurunkan scope-nya ke child key dengan batas waktu hidup yang lebih sempit. Saya memperbaiki provisioning organisasi agar \texttt{organization\_id} yang belum ada di tabel organisasi tidak menyebabkan kegagalan (PR \#6).
    \item \textbf{Registration role provisioning}: Memperbaiki provisioning role pada registrasi pengguna PostgreSQL dan menambahkan klarifikasi pembuatan tim dengan peran user pada Stack Auth.
    \item \textbf{DPoP Keycloak Tutorial}: Menulis tutorial end-to-end demonstrasi proof-of-possession (DPoP) dengan Keycloak, termasuk penanganan error yang jelas saat koneksi gagal, penghapusan kunci yang tidak sengaja ter-commit, serta penyesuaian path \texttt{pyproject.toml} (PR \#7).
\end{itemize}

\subsubsection{Deployment Digital Employee dengan Helm}

Saya berpartisipasi dalam proses deployment aplikasi Digital Employee ke lingkungan dev dan prod melalui pull request deploy yang menggerakkan pipeline CI/CD. Kontribusi teknis di sisi konfigurasi meliputi:

\begin{itemize}
    \item Menyelaraskan template ConfigMap dan Secret Helm dengan nilai dev, termasuk membalikkan konfigurasi placeholder GLC yang salah.
    \item Menata ulang \texttt{.env.example} dengan seksi Legacy untuk variabel MCP yang masih dipakai, dan menghapus default variabel \texttt{GL\_CONNECTORS\_API\_KEY}/\texttt{TOKEN} yang tidak lagi diperlukan setelah migrasi GLConnectors.
\end{itemize}

Dengan pemisahan antara kode aplikasi dan nilai konfigurasi, perubahan variabel lingkungan tidak memerlukan perubahan kode maupun rebuild image.

\subsection{Refleksi}

Perbaikan atomicitas \texttt{create\_user\_with\_password} mengajarkan bahwa solusi yang tampak cukup (rollback \texttt{try/except}) belum tentu benar secara arsitektural. Memindahkan tanggung jawab atomicitas ke penyedia data, tempat transaksi sebenarnya terjadi, menghasilkan kode yang lebih sederhana dan perilaku yang benar, sesuai prinsip separation of concerns. Pengalaman ini juga menunjukkan pentingnya menulis test untuk skenario kegagalan parsial, bukan hanya jalur sukses.

Dari pekerjaan di cookbook, saya belajar bahwa contoh integrasi sama pentingnya dengan SDK itu sendiri: dokumentasi yang jelas tentang provisioning role dan batasan waktu API key menyelamatkan pengguna dari kegagalan yang sulit didiagnosis.

\clearpage
